\documentclass[leqno]{article}
% Shared mark scheme preamble for resources.danieljsmith.org: the pack
% preamble, unchanged, then the scheme's own layer. A scheme's .tex is
% rendered from its YAML (scripts/README.md, Mark schemes), never edited.
\input{\CurrentFilePath/pack_preamble.tex}

% The page: equal side margins, so the table runs the full width of the
% header rule. The footers' right edge follows the text's.
\geometry{left=0.8in, right=0.8in}
\makeatletter
\renewcommand{\djs@footR}[1]{%
  \put(\dimexpr1in+\hoffset+\oddsidemargin+\textwidth\relax,
    -\dimexpr\paperheight-0.32in\relax){\makebox[0pt][r]{#1}}}
\makeatother

\usepackage[export]{adjustbox}
\usepackage{tabularray}
\definecolor{msGrid}{RGB}{176, 172, 164}
\definecolor{msTint}{RGB}{245, 243, 239}

% The header on every page, the front page included, and the booklet label
% under the contents-page title. \djsQuestionPageStyle stays empty.
\newcommand{\djsSchemeHeader}[1]{%
  \djsPackHeader[markscheme]{\textit{Mark scheme} - #1}%
  \renewcommand{\djsPackTitle}{#1}%
  \renewcommand{\djsBookletLabel}{Mark scheme}%
}

% The front page: the contents heading, spaced as tocloft spaces the packs'
% "Questions", then the contents table. A question's name and page number
% both link to its first page.
\newcommand{\msContentsHeading}{%
  \vskip\cftbeforetoctitleskip
  {\noindent\cfttoctitlefont\contentsname\par}%
  \vskip\cftaftertoctitleskip}
\NewTblrEnviron{mscontents}
\SetTblrInner[mscontents]{
  colspec = {Q[wd=26mm, l] X[l] Q[wd=16mm, c] Q[wd=14mm, r]},
  width = \linewidth,
  rowsep = 4pt, colsep = 6pt,
  hlines = {0.4pt, msGrid}, vlines = {0.4pt, msGrid},
  hline{2} = {2.2pt, scarlet},
  row{1} = {font=\small\bfseries},
  row{Z} = {font=\bfseries},
  column{2} = {font=\small},
}
\newcommand{\msContentsQuestion}[1]{\hyperref[q:#1]{Question #1}}
\newcommand{\msContentsPage}[1]{\hyperref[q:#1]{\pageref*{q:#1}}}
\newcommand{\msSep}{\hspace{0.9em}}

% The conventions pages' layout: headings in the documents' own bold, tables
% ruled and padded as the scheme's, roomy lists.
\newcommand{\msHead}[1]{%
  \par\vspace{18pt}%
  {\noindent\bfseries #1\par}%
  \vspace{6pt}}
\NewTblrEnviron{mstblr}
\SetTblrInner[mstblr]{rowsep = 6pt, colsep = 6pt,
  hlines = {0.4pt, msGrid}, vlines = {0.4pt, msGrid}}
\newenvironment{msList}
  {\begin{itemize}[leftmargin=1.2em, itemsep=5pt, topsep=3pt, parsep=0pt]}
  {\end{itemize}}

% A question's first page: the link target and the PDF bookmark, as
% \questionitem gives the packs.
\newcommand{\msQnum}{}
\newcommand{\msQuestion}[1]{%
  \clearpage
  \renewcommand{\msQnum}{#1}%
  \phantomsection\label{q:#1}%
  \addcontentsline{toc}{subsection}{\protect\djsTocEntry{#1}}}

% The scheme's table: Question 16mm, Scheme the rest, Marks 16mm, Notes
% 44mm; every cell ruled in msGrid; the scarlet rule under the head row. A
% long table has no caption or "continued" foot; a continuation page opens
% with a muted line above the repeated head row. The rules across a part's
% rows, the alternative's tint and the total row are the rendered table's.
\NewTblrTheme{ms}{%
  \DefTblrTemplate{firsthead}{default}{}%
  \DefTblrTemplate{middlehead,lasthead}{default}{%
    {\small\itshape\color{djsMuted}Question \msQnum\ continued\par}%
    \vspace{6pt}}%
  \DefTblrTemplate{firstfoot,middlefoot,lastfoot}{default}{}%
}
\NewTblrEnviron{msscheme}
\SetTblrOuter[msscheme]{long, theme=ms, entry=none, label=none}
\SetTblrInner[msscheme]{
  colspec = {Q[wd=16mm, l] X[l] Q[wd=16mm, c] Q[wd=44mm, l]},
  width = \linewidth,
  rowsep = 6pt, colsep = 6pt,
  rowhead = 1,
  vlines = {0.4pt, msGrid},
  hline{1} = {0.4pt, msGrid}, hline{2} = {2.2pt, scarlet},
  hline{Z} = {0.4pt, msGrid},
  column{1} = {font=\bfseries}, column{3} = {font=\bfseries},
  column{4} = {font=\small},
  row{1} = {rowsep=4pt, font=\small\bfseries},
}

% Two marks on a row: each note starts with its code.
\newcommand{\msNote}[2]{\textbf{#1}\enspace #2}
% A full-row figure: the worked solution's own picture, scaled uniformly and
% only down, to the table less the Question column. The picture is drawn
% with the line width it has in the solutions document, at its list depth
% there (the optional argument, 1 by default: a part's working), so one
% sized by \linewidth keeps its size; tabularray sets the cell's \linewidth
% to the cell's measured width, which would shrink it at every pass. The
% widths are measured once, in a solution box at the packs' text width
% (pack_preamble.tex's margins, 0.8in and 1in).
\newlength{\msSpanW}
\makeatletter
\newcommand{\ms@keepLW}[1]{%
  \expandafter\xdef\csname ms@solnLW@#1\endcsname{\the\linewidth}}
\AtBeginDocument{%
  \setlength{\msSpanW}{\dimexpr\textwidth-16mm-28pt\relax}%
  \setbox\z@\vbox{%
    \begin{minipage}{\dimexpr\paperwidth-1.8in\relax}
      \begin{djsSolution}
        \ms@keepLW{0}%
        \begin{enumerate}\item \ms@keepLW{1}%
        \begin{enumerate}\item \ms@keepLW{2}%
        \begin{enumerate}\item \ms@keepLW{3}%
        \end{enumerate}\end{enumerate}\end{enumerate}
      \end{djsSolution}
    \end{minipage}}}
\newcommand{\msFigure}[2][1]{%
  \adjustbox{max width=\msSpanW, max totalheight=0.6\textheight, valign=t}{%
    \hsize=\csname ms@solnLW@#1\endcsname
    \linewidth=\hsize
    #2}}
\makeatother
% A figure mark's feature, in the Scheme column.
\newcommand{\msFeature}[1]{\textit{#1}}
% An alternative's title row.
\newcommand{\msAlternative}[1]{\textit{Alternative:} #1}
% Drafts only: a queried mark's code, the query's why under its note, and a
% flag in its part's place.
\newcommand{\msQueried}[2]{#1~\(\to\) #2?}
\newcommand{\msQuery}[1]{{\itshape\color{djsMuted}Query: #1}}
\newcommand{\msFlag}[2]{\textbf{#1:}\enspace #2}


\djsSchemeHeader{Proof by Induction - Series}
\djsSitePath{Core Pure}{Proof by Induction}

\djsLastUpdated{25 October 2026}
\begin{document}
\djsSetQuestionCount{10}
\djsAddMarks{1}{6}
\djsAddMarks{2}{6}
\djsAddMarks{3}{6}
\djsAddMarks{4}{6}
\djsAddMarks{5}{12}
\djsAddMarks{6}{8}
\djsAddMarks{7}{6}
\djsAddMarks{8}{6}
\djsAddMarks{9}{11}
\djsAddMarks{10}{6}
\vspace*{6pt}
\djsContentsTitle
\vspace{30pt}
\msContentsHeading
\begin{mscontents}{}
Question & Marks by part & Marks & Page \\
\msContentsQuestion{1} & & 6 & \msContentsPage{1} \\
\msContentsQuestion{2} & & 6 & \msContentsPage{2} \\
\msContentsQuestion{3} & & 6 & \msContentsPage{3} \\
\msContentsQuestion{4} & & 6 & \msContentsPage{4} \\
\msContentsQuestion{5} & (a)\,6\msSep (b)\,1\msSep (c)\,2\msSep (d)\,3 & 12 & \msContentsPage{5} \\
\msContentsQuestion{6} & (a)\,6\msSep (b)\,2 & 8 & \msContentsPage{6} \\
\msContentsQuestion{7} & & 6 & \msContentsPage{7} \\
\msContentsQuestion{8} & & 6 & \msContentsPage{8} \\
\msContentsQuestion{9} & (a)\,6\msSep (b)\,5 & 11 & \msContentsPage{9} \\
\msContentsQuestion{10} & & 6 & \msContentsPage{10} \\
\SetCell[c=2]{r} Total & & 73 & \\
\end{mscontents}
\newpage
% The conventions pages, shared by every pack's mark scheme: one source of
% truth, \input after the front page, with no strand or pack logic. The keys
% on the first side, the rules on the second. The scheme layer defines
% \msHead, msList and mstblr.
\vspace*{6pt}%
{\noindent\cfttoctitlefont Reading the mark scheme\par}%
\vspace{10pt}%
These rules apply throughout, unless a note says otherwise.

\msHead{Marks}
Every mark is worth one, so \textbf{M1 A1} on a row is two marks. Only marks
coded ft follow through.\par
\vspace{10pt}%
\begin{mstblr}{
  colspec = {Q[wd=15mm, l] X[l]},
  width = \linewidth,
  hline{2} = {2.2pt, scarlet},
  column{1} = {font=\bfseries},
  row{1} = {font=\small\bfseries, rowsep=4pt},
}
Code & Meaning \\
M1 & Method: a correct method, applied to this question. A slip doesn't cost it if the method is clear. A formula must be used, not just quoted. \\
dM1 & Dependent method: as M1, but only if the method mark before it was earned, unless its note names another. \\
A1 & Accuracy: a correct answer or step. It needs the method mark it follows. \\
A1ft & As A1, following through from an earlier answer that the note names. \\
A1~AG & The last mark of a part whose answer is printed in the question: complete working, with no errors. \\
B1 & A correct value, fact or statement. It needs no method mark. \\
B1ft & As B1, following through from an earlier answer that the note names. \\
E1 & A written explanation, interpretation or conclusion. Concluding a proof by induction needs every earlier mark. \\
\end{mstblr}

\msHead{Abbreviations}
ft and AG qualify a code in the Marks column. The others appear in the
notes.\par
\vspace{10pt}%
\begin{mstblr}{
  colspec = {Q[wd=17mm, l] X[l] Q[wd=15mm, l] X[l]},
  width = \linewidth,
  column{1} = {font=\bfseries}, column{3} = {font=\bfseries},
}
ft & follow through & AG & answer given in the question \\
cao & correct answer only, with no follow-through & cso & correct solution only: no errors in the part \\
oe & or an equivalent form & awrt & any answer that rounds to this value \\
soi & seen, or implied by later working & isw & ignore working after a correct answer \\
SC & special case: credit outside the main method & their & the student's own earlier value \\
condone & a slip that doesn't cost the mark & [\;\,] & working the mark doesn't need \\
s.f. & significant figures & d.p. & decimal places \\
\end{mstblr}

\newpage
\msHead{Reading the scheme}
\begin{msList}
\item Each row is a line of working, with the marks it earns beside it.
  Where a row carries two marks, each note starts with its code.
\item \textbf{[\textit{n}]} ends each part with its total.
\item A shaded \textit{Alternative} is a second method for the same total.
  Mark an answer by the method it follows.
\item Any other valid method earns the equivalent marks, unless the question
  names a method or says \textit{hence}.
\item A sketch earns one B1 for each feature named below its figure.
\end{msList}

\msHead{Accuracy and answers}
\begin{msList}
\item Unless a question says otherwise, give answers to 3 significant
  figures, and angles in degrees to 1 decimal place. Angles in radians take
  3 significant figures.
\item A final answer with more figures is accepted if it rounds to the
  scheme's value; one with fewer is not. An answer spoiled by rounding too
  early loses its accuracy mark.
\item Where a question asks for an exact answer, a decimal loses the final
  accuracy mark. Otherwise either is accepted.
\item Wrong or missing units cost nothing unless the question asks for them.
\item A correct answer with no working earns full marks, unless the question
  needs the working: show that, prove, hence, a method it names, or an exact
  answer. A wrong answer with no working earns nothing. Calculator answers
  are treated the same way.
\item A correct answer, once reached, isn't spoiled by later working (isw).
\end{msList}

\msHead{Value of $\boldsymbol{g}$}
\begin{msList}
\item Take $g = 9.8\,\mathrm{m\,s^{-2}}$ unless a question states a value,
  which is then the only one accepted. Answers that use $g = 9.8$ may be
  given to 2 or 3 significant figures.
\item Where a question doesn't state $g$, an answer worked with $g = 10$
  throughout earns the same marks, keeping any value printed in the
  question. The worked solutions list the final answers with $g = 10$.
\end{msList}


\msQuestion{1}
\begin{msscheme}{
  hline{3-7} = {2-Z}{0.4pt, msGrid},
  hline{8} = {0.4pt, msGrid},
  row{8} = {font=\bfseries},
}
Question & Scheme & Marks & Notes \\
1 & {\(n=1\): LHS \(= 1\times 2 = 2\), RHS \(= 2+(1-1)2^{2} = 2\)} & B1 & {Both sides evaluated at \(n=1\) and shown equal to 2.} \\
& {Assume \(\sum_{r=1}^{k} r\,2^r = 2+(k-1)2^{k+1}\)\\[4pt] \(\sum_{r=1}^{k+1} r\,2^r = \sum_{r=1}^{k} r\,2^r + (k+1)2^{k+1}\)} & M1 & {States the assumption for \(n=k\) and adds a \((k+1)\)th term of the form \((k+1)2^{\ldots}\) to the sum to \(k\). Condone a slip in the power.} \\
& {\(= 2+(k-1)2^{k+1}+(k+1)2^{k+1}\)\\[4pt] \(= 2+\bigl((k-1)+(k+1)\bigr)2^{k+1}\)} & M1 & {Replaces the sum to \(k\) by the hypothesis and collects the terms in \(2^{k+1}\), or writes both with a common power of 2. Condone sign slips.} \\
& {\(= 2+2k\,2^{k+1}\)} & A1 & {oe, such as \(2+k\,2^{k+2}\).} \\
& {\(= 2+k\,2^{k+2} = 2+\bigl((k+1)-1\bigr)2^{(k+1)+1}\)} & A1 & {cso. \(2+\bigl((k+1)-1\bigr)2^{(k+1)+1}\) seen, or \(2+k\,2^{k+2}\) explicitly compared with the result at \(n=k+1\).} \\
& {True for \(n=1\); true for \(n=k\) \(\Rightarrow\) true for \(n=k+1\); so true for all \(n \geq 1\)} & E1 & {States all three points: true for \(n=1\), true for \(n=k\) implies true for \(n=k+1\), so true for all \(n \geq 1\). One sentence accepted.} \\
\nopagebreak \SetCell[c=2]{r} Total for Question 1 & & 6 & \\
\end{msscheme}

\msQuestion{2}
\begin{msscheme}{
  hline{3-6} = {2-Z}{0.4pt, msGrid},
  hline{7} = {0.4pt, msGrid},
  row{7} = {font=\bfseries},
}
Question & Scheme & Marks & Notes \\
2 & {\(n=1\): LHS \(= 2 + 2\times 3 = 8\)\\[4pt] \(\phantom{n=1:}\) RHS \(= 2^2 + 3^2 - 5 = 8\)} & B1 & {Both sides evaluated to 8; a bare 8 = 8 is not enough.} \\
& {Assumes true for \(n=k\):\\[4pt] \(\sum_{r=1}^{k+1} = 2^{k+1}+3^{k+1}-5\)\\[4pt] \(\phantom{\sum_{r=1}^{k+1}} + 2^{k+1}+2\times 3^{k+1}\)} & M1 & {States the assumption for \(n=k\) and adds a \((k+1)\)th term to its RHS, or shows RHS\((k+1)\) \(-\) RHS\((k)\) is that term. Condone a slip in the term.} \\
& {\(\phantom{\sum_{r=1}^{k+1}} = 2\times 2^{k+1} + 3\times 3^{k+1} - 5\)} & M1 A1 & {\msNote{M1}{Collects like powers of 2 and of 3 from their sum, e.g. \(2^{k+1}+2^{k+1}\) to \(2\times 2^{k+1}\).}\\ \msNote{A1}{Correct, oe.}} \\
& {\(\phantom{\sum_{r=1}^{k+1}} = 2^{k+2}+3^{k+2}-5\)} & A1 & {cso; reached from their sum, or \(2^{k+2}+3^{k+2}-5\) shown equal to their expression; A0 if asserted.} \\
& {True for \(n=1\); true for \(n=k\)\\[4pt] \(\Rightarrow\) true for \(n=k+1\);\\[4pt] so true for all \(n \geq 1\)} & E1 & {States all three: true for \(n=1\), true for \(n=k\) implies true for \(n=k+1\), hence true for all \(n \geq 1\).} \\
\nopagebreak \SetCell[c=2]{r} Total for Question 2 & & 6 & \\
\end{msscheme}

\msQuestion{3}
\begin{msscheme}{
  hline{3-6} = {2-Z}{0.4pt, msGrid},
  hline{7} = {0.4pt, msGrid},
  row{7} = {font=\bfseries},
}
Question & Scheme & Marks & Notes \\
3 & {\(n=1\): LHS \(= 1(2)(3) = 6\)\\[4pt] \(\phantom{n=1:}\) RHS \(= \tfrac14(1)(2)(3)(4) = 6\)} & B1 & {Both sides evaluated at \(n=1\); ``true for \(n=1\)'' alone is not enough.} \\
& {Assumes true for \(n=k\):\\[4pt] \(\sum_{r=1}^{k+1} r(r+1)(r+2)\)\\[4pt] \(\phantom{A} = \tfrac14k(k+1)(k+2)(k+3) + (k+1)(k+2)(k+3)\)} & M1 & {States the assumption for \(n=k\), in words or symbols, and adds the \((k+1)\)th term \((k+1)(k+2)(k+3)\) to the sum to \(k\).} \\
& {\(\phantom{A} = (k+1)(k+2)(k+3)\left(\tfrac{k}{4}+1\right)\)} & M1 A1 & {\msNote{M1}{Uses the hypothesis and takes out a common factor of at least two of \(k+1\), \(k+2\), \(k+3\), or expands both terms into one quartic. Condone slips in the remaining bracket.}\\ \msNote{A1}{Correct factorised expression oe, or \(\frac14(k^4+10k^3+35k^2+50k+24)\).}} \\
& {\(\phantom{A} = \tfrac14(k+1)(k+2)(k+3)(k+4)\)} & A1 & {cso, oe in terms of \(k+1\). A quartic must be shown equal to the target with \(n=k+1\).} \\
& {True for \(n=1\); true for \(n=k\) \(\Rightarrow\) true for \(n=k+1\); so true for all positive integers \(n\)} & E1 & {States all three points: true for \(n=1\), true for \(n=k\) implies true for \(n=k+1\), hence true for all positive integers \(n\).} \\
\nopagebreak \SetCell[c=2]{r} Total for Question 3 & & 6 & \\
\end{msscheme}

\msQuestion{4}
\begin{msscheme}{
  hline{3-7} = {2-Z}{0.4pt, msGrid},
  hline{8} = {0.4pt, msGrid},
  row{8} = {font=\bfseries},
}
Question & Scheme & Marks & Notes \\
4 & {\(n=1\): LHS \(= 1\times 3^0 = 1\)\\[4pt] RHS \(= \frac{1+(2-1)3^1}{4} = 1\)} & B1 & {Both sides shown to equal 1; a bare statement that it is true for \(n=1\) is not enough.} \\
& {Assume true for \(n=k\):\\[4pt] \(\sum_{r=1}^{k+1} r\,3^{r-1}\)\\[4pt] \(\phantom{A} = \sum_{r=1}^{k} r\,3^{r-1} + (k+1)3^k\)} & M1 & {States the assumption for \(n=k\) and adds a \((k+1)\)th term of the form \((k+1)3^{p}\) to the sum to \(k\). Condone a slip in the power.} \\
& {\(\phantom{A} = \frac{1+(2k-1)3^k + 4(k+1)3^k}{4}\)} & M1 & {Replaces the sum to \(k\) by the hypothesis and writes the result as a single fraction over 4 or a multiple of 4. Condone sign slips.} \\
& {\(\phantom{A} = \frac{1+(6k+3)3^k}{4}\)} & A1 & {Correct numerator collected as one multiple of \(3^k\), oe.} \\
& {\(\phantom{A} = \frac{1+(2k+1)3^{k+1}}{4}\)} & A1 & {cso. Reached by factorising \(3^{k+1}\), or by simplifying the RHS at \(n=k+1\) to the same expression as their sum.} \\
& {True for \(n=1\); true for \(n=k\)\\[4pt] \(\Rightarrow\) true for \(n=k+1\)\\[4pt] \(\Rightarrow\) true for all \(n\in\mathbb{Z}^+\)} & E1 & {States true for \(n=1\), true for \(n=k\) implies true for \(n=k+1\), and hence true for all positive integers \(n\); any order.} \\
\nopagebreak \SetCell[c=2]{r} Total for Question 4 & & 6 & \\
\end{msscheme}

\msQuestion{5}
\begin{msscheme}{
  hline{3-8,10,12-13,15-16} = {2-Z}{0.4pt, msGrid},
  hline{9,11,14,17} = {0.4pt, msGrid},
  row{17} = {font=\bfseries},
}
Question & Scheme & Marks & Notes \\
5(a) & {\(n=1\): LHS \(= 1^3 = 1\), RHS \(= \frac14(1)^2(2)^2 = 1\)} & B1 & {Both sides shown equal to 1.} \\
& {Assume \(\sum_{r=1}^{k} r^3 = \frac14k^2(k+1)^2\)\\[4pt] \(\sum_{r=1}^{k+1} r^3 = \sum_{r=1}^{k} r^3 + (k+1)^3\)} & M1 & {Assumes the result for \(n=k\) and adds \((k+1)^3\) to the sum to \(k\).} \\
& {\(\phantom{\sum} = \frac14k^2(k+1)^2 + (k+1)^3\)\\[4pt] \(\phantom{\sum} = (k+1)^2(\ldots)\)} & M1 & {Substitutes the hypothesis and takes out a factor \((k+1)^2\) or \(\frac14(k+1)^2\); or expands to a quartic for comparison with \(\frac14(k+1)^2(k+2)^2\) expanded.} \\
& {\(\phantom{\sum} = (k+1)^2\left(\frac{k^2}{4}+k+1\right)\)} & A1 & {Correct expression, oe such as \(\frac14(k^4+6k^3+13k^2+12k+4)\).} \\
& {\(\phantom{\sum} = \frac14(k+1)^2(k+2)^2\)} & A1 & {Reaches \(\frac14(k+1)^2(k+2)^2\), or shows it equals their correct quartic, without error.} \\
& {True for \(n=1\); true for \(n=k \Rightarrow\) true for \(n=k+1\)\\[4pt] so true for all integers \(n \ge 1\)} & E1 & {States true for \(n=1\), true for \(n=k\) implies true for \(n=k+1\), hence true for all \(n\ge1\); may be split across the proof.} \\
\nopagebreak & & [6] & \\
5(b) & {\(\frac14(2n)^2(2n+1)^2 = n^2(2n+1)^2\)} & B1 & {Or any equivalent factorised form. Not an expanded polynomial.} \\
\nopagebreak & & [1] & \\
5(c) & {\(\sum_{r=1}^{n}(2r)^3 = 8\sum_{r=1}^{n} r^3 = 8\times\frac14n^2(n+1)^2\)} & M1 & {Takes out the factor \(8\) or \(2^3\) and applies the formula of part (a) with upper limit \(n\).} \\
& {\(\phantom{\sum} = 2n^2(n+1)^2\)} & A1 & {Or any equivalent factorised form.} \\
\nopagebreak & & [2] & \\
5(d) & {First \(2n\) cubes \(-\) even cubes:\\[4pt] \(1^3+3^3+\cdots+(2n-1)^3\)\\[4pt] \(\phantom{1} = n^2(2n+1)^2 - 2n^2(n+1)^2\)} & M1 A1 & {\msNote{M1}{Subtracts their (c) from their (b); or expands \(\sum(2r-1)^3\) and applies standard summation formulae.}\\ \msNote{A1}{Correct difference, or correct unsimplified sum by the other route.}} \\
& {\(\phantom{1} = n^2(2n^2-1)\), \(a=1\), \(b=2\)} & A1 & {cso. \(a=1\), \(b=2\) stated, or \(n^2(2n^2-1)\) reached.} \\
\nopagebreak & & [3] & \\
\nopagebreak \SetCell[c=2]{r} Total for Question 5 & & 12 & \\
\end{msscheme}

\msQuestion{6}
\begin{msscheme}{
  hline{3-7,9-10} = {2-Z}{0.4pt, msGrid},
  hline{8,11} = {0.4pt, msGrid},
  row{11} = {font=\bfseries},
}
Question & Scheme & Marks & Notes \\
6(a) & {\(n=1\): LHS \(= \frac{1}{1\times 3} = \frac13\), RHS \(= \frac{1}{2+1} = \frac13\)} & B1 & {Both sides evaluated as \(\frac13\); a bare statement that it is true for \(n=1\) is not enough.} \\
& {Assume true for \(n=k\):\\[4pt] \(\sum_{r=1}^{k+1} = \frac{k}{2k+1} + \frac{1}{(2k+1)(2k+3)}\)} & M1 & {States the assumption for \(n=k\) and adds the \((k+1)\)th term to the sum of \(k\) terms. The term need not be simplified.} \\
& {\(\phantom{\sum} = \frac{k(2k+3)+1}{(2k+1)(2k+3)}\)} & M1 A1 & {\msNote{M1}{Replaces the sum of \(k\) terms by \(\frac{k}{2k+1}\) and combines over any common denominator. Condone slips in the numerator.}\\ \msNote{A1}{Any correct single fraction, unsimplified.}} \\
& {\(\phantom{\sum} = \frac{(2k+1)(k+1)}{(2k+1)(2k+3)} = \frac{k+1}{2(k+1)+1}\)} & A1 & {cso. Correct factorisation reaching \(\frac{k+1}{2k+3}\) or \(\frac{k+1}{2(k+1)+1}\).} \\
& {True for \(n=1\); true for \(n=k\) \(\Rightarrow\) true for \(n=k+1\); so true for all positive integers \(n\)} & E1 & {States all three points: true for \(n=1\), true for \(n=k\) implies true for \(n=k+1\), hence true for all positive integers \(n\).} \\
\nopagebreak & & [6] & \\
6(b) & {\(\sum_{r=n+1}^{2n} = \frac{2n}{4n+1} - \frac{n}{2n+1}\)\\[4pt] \(\phantom{\sum} = \frac{2n(2n+1) - n(4n+1)}{(4n+1)(2n+1)}\)} & M1 & {Replaces \(n\) by \(2n\) in the result of (a), subtracts the sum of \(n\) terms and combines over a common denominator. Condone slips in the numerator.} \\
& {\(\phantom{\sum} = \frac{n}{(4n+1)(2n+1)}\)\\[4pt] \(a=4,\ b=1,\ c=2,\ d=1\)} & A1 & {cso. Factors in either order, so \(a=2,\ b=1,\ c=4,\ d=1\) also scores.} \\
\nopagebreak & & [2] & \\
\nopagebreak \SetCell[c=2]{r} Total for Question 6 & & 8 & \\
\end{msscheme}

\msQuestion{7}
\begin{msscheme}{
  hline{3-6} = {2-Z}{0.4pt, msGrid},
  hline{7} = {0.4pt, msGrid},
  row{7} = {font=\bfseries},
}
Question & Scheme & Marks & Notes \\
7 & {\(n=1\): LHS \(= \frac12\), RHS \(= 2 - \frac32 = \frac12\)} & B1 & {Both sides shown equal to \(\frac12\); not for one side alone.} \\
& {Assume \(\sum_{r=1}^{k}\frac{r}{2^r} = 2 - \frac{k+2}{2^k}\)\\[4pt] \(\sum_{r=1}^{k+1}\frac{r}{2^r} = 2 - \frac{k+2}{2^k} + \frac{k+1}{2^{k+1}}\)} & M1 & {States or uses the result for \(n=k\) and adds a \((k+1)\)th term to it. Condone a slip in the term.} \\
& {\(\phantom{\sum_{r=1}^{k+1}\frac{r}{2^r}} = 2 - \frac{2(k+2)-(k+1)}{2^{k+1}}\)} & M1 A1 & {\msNote{M1}{Combines their expression over a common denominator such as \(2^{k+1}\). Condone sign slips.}\\ \msNote{A1}{Correct single fraction oe; need not be simplified.}} \\
& {\(\phantom{\sum_{r=1}^{k+1}\frac{r}{2^r}} = 2 - \frac{(k+1)+2}{2^{k+1}}\)} & A1 & {cso. \(2-\frac{k+3}{2^{k+1}}\) scores only if linked to the result for \(n=k+1\), stated as target or conclusion.} \\
& {True for \(n=1\); true for \(n=k \Rightarrow\) true for \(n=k+1\)\\[4pt] so true for all \(n \ge 1\)} & E1 & {States true for \(n=1\), truth for \(n=k\) implies truth for \(n=k+1\), hence true for all \(n \ge 1\). Needs the assumption for \(n=k\) stated somewhere.} \\
\nopagebreak \SetCell[c=2]{r} Total for Question 7 & & 6 & \\
\end{msscheme}

\msQuestion{8}
\begin{msscheme}{
  hline{3-7} = {2-Z}{0.4pt, msGrid},
  hline{8} = {0.4pt, msGrid},
  row{8} = {font=\bfseries},
}
Question & Scheme & Marks & Notes \\
8 & {\(n=1\): LHS \(= \frac{1}{1\times2\times3} = \frac16\)\\[4pt] RHS \(= \frac{1\times4}{4\times2\times3} = \frac16\)} & B1 & {Both sides evaluated as \(\frac16\). Not for a bare ``true for \(n=1\)''.} \\
& {Assume \(\sum_{r=1}^{k} \frac{1}{r(r+1)(r+2)} = \frac{k(k+3)}{4(k+1)(k+2)}\)\\[4pt] \(\sum_{r=1}^{k+1} \frac{1}{r(r+1)(r+2)}\)\\[4pt] \(\phantom{A} = \sum_{r=1}^{k} \frac{1}{r(r+1)(r+2)} + \frac{1}{(k+1)(k+2)(k+3)}\)} & M1 & {States the assumption for \(n=k\) and adds the \((k+1)\)th term to the sum to \(k\) terms. Condone a slip in the added term.} \\
& {\(\phantom{A} = \frac{k(k+3)}{4(k+1)(k+2)} + \frac{1}{(k+1)(k+2)(k+3)}\)\\[4pt] \(\phantom{A} = \frac{\ldots}{4(k+1)(k+2)(k+3)}\)} & M1 & {Replaces the sum to \(k\) terms by the hypothesis and combines with their \((k+1)\)th term over a common denominator. Condone numerator slips.} \\
& {\(\phantom{A} = \frac{k(k+3)^2+4}{4(k+1)(k+2)(k+3)}\)} & A1 & {Any correct single fraction, unsimplified or with the numerator expanded.} \\
& {\(\phantom{A} = \frac{(k+1)^2(k+4)}{4(k+1)(k+2)(k+3)}\)\\[4pt] \(\phantom{A} = \frac{(k+1)(k+4)}{4(k+2)(k+3)}\)} & A1 & {cso, oe in terms of \(k+1\): numerator factorised, or the \(n=k+1\) result shown equal to their fraction; A0 if asserted.} \\
& {True for \(n=1\); true for \(n=k \Rightarrow\) true for \(n=k+1\)\\[4pt] \(\Rightarrow\) true for all \(n \in \mathbb{Z}^+\)} & E1 & {Concludes with all three points: true for \(n=1\); if true for \(n=k\) then true for \(n=k+1\); so true for all \(n \in \mathbb{Z}^+\).} \\
\nopagebreak \SetCell[c=2]{r} Total for Question 8 & & 6 & \\
\end{msscheme}

\msQuestion{9}
\begin{msscheme}{
  hline{3-8,10-14} = {2-Z}{0.4pt, msGrid},
  hline{9,15} = {0.4pt, msGrid},
  row{15} = {font=\bfseries},
}
Question & Scheme & Marks & Notes \\
9(a) & {\(n=1\): LHS \(= 1^3 = 1\), RHS \(= 1^2(2-1) = 1\)} & B1 & {Both sides shown equal to 1; a bare ``true for \(n=1\)'' is not enough.} \\
& {Assumes \(\sum_{r=1}^{k}(2r-1)^3 = k^2(2k^2-1)\)\\[4pt] \(\sum_{r=1}^{k+1}(2r-1)^3 = \sum_{r=1}^{k}(2r-1)^3 + (2k+1)^3\)} & M1 & {Assumes the result for \(n=k\) and adds the \((k+1)\)th term, \((2k+1)^3\) or \((2(k+1)-1)^3\), to the sum to \(k\).} \\
& {\(\phantom{\sum} = k^2(2k^2-1) + (2k+1)^3 = \ldots\)} & M1 & {Substitutes \(k^2(2k^2-1)\) for the sum to \(k\) and expands or factorises towards the \(n=k+1\) form. Condone slips in expanding the cube.} \\
& {\(\phantom{\sum} = 2k^4+8k^3+11k^2+6k+1\)} & A1 & {oe, such as a correct factorised form.} \\
& {\((k+1)^2(2(k+1)^2-1) = (k^2+2k+1)(2k^2+4k+1)\)\\[4pt] \(\phantom{(k+1)^2} = 2k^4+8k^3+11k^2+6k+1\)} & A1 & {Shows the quartic equals \((k+1)^2(2(k+1)^2-1)\), by expanding this or factorising the quartic, without error.} \\
& {True for \(n=1\); true for \(n=k \Rightarrow\) true for \(n=k+1\)\\[4pt] \(\Rightarrow\) true for all positive integers \(n\)} & E1 & {States true for \(n=1\), true for \(n=k\) implies true for \(n=k+1\), and so true for all positive integers \(n\); may be spread through the working.} \\
\nopagebreak & & [6] & \\
9(b) & {\(\sum_{r=1}^{2n}(2r-1)^3 = (2n)^2\left(2(2n)^2-1\right)\)} & M1 & {Replaces \(n\) by \(2n\) throughout the formula from (a). Need not be simplified.} \\
& {\(\phantom{\sum} = 4n^2(8n^2-1)\)} & A1 & {oe, such as \(32n^4-4n^2\).} \\
& {\(\their{4n^2(8n^2-1)} = kn^2(2n^2-1)\)} & M1 & {Sets their sum to \(2n\), obtained from (a)'s formula, equal to \(k\) times \(n^2(2n^2-1)\).} \\
& {\(32n^2-4 = 2kn^2-k \Rightarrow (32-2k)n^2 = 4-k\)} & dM1 & {Divides by \(n^2\) and collects the \(n^2\) terms to reach the form \(n^2(\ldots) = \ldots\). Condone sign slips.} \\
& {\(n^2 = \frac{k-4}{2(k-16)}\)} & A1~AG & {Correct substituted equation leading without error to the printed form; combined steps accepted.} \\
\nopagebreak & & [5] & \\
\nopagebreak \SetCell[c=2]{r} Total for Question 9 & & 11 & \\
\end{msscheme}

\msQuestion{10}
\begin{msscheme}{
  hline{3-7} = {2-Z}{0.4pt, msGrid},
  hline{8} = {0.4pt, msGrid},
  row{8} = {font=\bfseries},
}
Question & Scheme & Marks & Notes \\
10 & {\(n=1\): LHS \(= \log\left(1-\frac14\right) = \log\frac34\)\\[4pt] RHS \(= \log\left(\frac{2!\,3!}{2^4}\right) = \log\frac34\)} & B1 & {Both sides shown equal to \(\log\frac34\), oe; the RHS must be evaluated, not just stated.} \\
& {Assumes true for \(n=k\):\\[4pt] \(\sum_{r=1}^{k+1}\ldots = \log\left(\frac{(2k)!(2k+1)!}{2^{4k}(k!)^4}\right)\)\\[4pt] \(\phantom{\sum_{r=1}^{k+1}\ldots} + \log\left(1-\frac{1}{4(k+1)^2}\right)\)} & M1 & {States the assumption for \(n=k\) and adds the \((k+1)\)th term to the assumed sum.} \\
& {\(= \log\left(\frac{(2k)!(2k+1)!}{2^{4k}(k!)^4}\times\frac{4(k+1)^2-1}{4(k+1)^2}\right)\)} & M1 & {Combines the two logs into one and writes the extra term as a single fraction. Condone slips in expanding.} \\
& {\(= \log\left(\frac{(2k)!(2k+1)!(2k+1)(2k+3)}{2^{4k}(k!)^4(2k+2)^2}\right)\)} & A1 & {Correct single log with the extra term as \(\frac{(2k+1)(2k+3)}{(2k+2)^2}\), oe such as \(4(k+1)^2\) in the denominator.} \\
& {\(= \log\left(\frac{(2k+2)!(2k+3)!}{2^{4k}(k!)^4(2k+2)^4}\right)\)\\[4pt] \(= \log\left(\frac{(2k+2)!(2k+3)!}{2^{4(k+1)}((k+1)!)^4}\right)\)} & A1 & {cso. Factorial identities for \((2k+2)!\), \((2k+3)!\) and \((k+1)!\) used without error to reach the \(n=k+1\) form; working from both ends to a common expression accepted.} \\
& {True for \(n=1\); true for \(n=k\) \(\Rightarrow\) true for \(n=k+1\); so true for all positive integers \(n\)} & E1 & {States true for \(n=1\), that truth for \(n=k\) implies truth for \(n=k+1\), and hence true for all positive integers \(n\).} \\
\nopagebreak \SetCell[c=2]{r} Total for Question 10 & & 6 & \\
\end{msscheme}

\end{document}
